
\documentclass{article}
\usepackage{colortbl}
\usepackage{makecell}
\usepackage{multirow}
\usepackage{supertabular}

\begin{document}

\newcounter{utterance}

\centering \large Interaction Transcript for game `codenames', experiment `frequency\_high', episode 7 with reflexion{-}Qwen3.5{-}4B.
\vspace{24pt}

{ \footnotesize  \setcounter{utterance}{1}
\setlength{\tabcolsep}{0pt}
\begin{supertabular}{c@{$\;$}|p{.15\linewidth}@{}p{.15\linewidth}p{.15\linewidth}p{.15\linewidth}p{.15\linewidth}p{.15\linewidth}}
    \# & $\;$A & \multicolumn{4}{c}{Game Master} & $\;\:$B\\
    \hline

    \theutterance \stepcounter{utterance}  
    & & & \multicolumn{2}{p{0.3\linewidth}}{
        \cellcolor[rgb]{0.9,0.9,0.9}{
            \makecell[{{p{\linewidth}}}]{
                \texttt{\tiny{[GM$|$GM]}}
                {'hidden': {'team': ['way', 'cause', 'report', 'mind', 'kid', 'army', 'middle', 'president'], 'innocent': ['story', 'coffee', 'light', 'evening', 'company', 'ship'], 'opponent': ['law', 'doctor', 'world', 'city', 'record', 'gold', 'fool'], 'assassin': ['hospital']}, 'revealed': {'team': {'team': [], 'innocent': [], 'opponent': [], 'assassin': []}, 'opponent': {'team': [], 'innocent': [], 'opponent': [], 'assassin': []}}}
            }
        }
    }
    & & \\ \\

    \theutterance \stepcounter{utterance}  
    & & \multicolumn{4}{p{0.6\linewidth}}{
        \cellcolor[rgb]{0.9,0.9,0.9}{
            \makecell[{{p{\linewidth}}}]{
                \texttt{\tiny{[P1$\langle$GM]}}
                \texttt{You have attempted this same task before and did not succeed. These are the plans you wrote after those attempts.} \\
\texttt{Trial \#1: New plan: The last attempt was aborted because the clue word "person" was identical to one of the target words, which is explicitly forbidden by the rules ("it cannot be one of the words in the lists"). Next time, I will select a clue word that is semantically related to at least two of my target words (e.g., "head" for mind, person, president) but is not present in the target, opponent, or distractor lists, ensuring the format strictly follows "CLUE: <WORD>" and "TARGETS: <WORDS>" with no extra text.} \\
\\ 
\texttt{Trial \#2: New plan: The last attempt was aborted because the clue word "head" is semantically related to "person" and "president," but the response format was likely rejected due to missing the required comma{-}separated list of targets or containing extra text. Next time, I will select a single{-}word clue that is not in any list (e.g., "face" for person, president, way) and strictly output only the two required lines: "CLUE: <WORD>" followed immediately by "TARGETS: <word1>, <word2>, ..." with no additional explanation or formatting.} \\
\\ 
\texttt{Trial \#3: New plan: The last attempt was aborted because the clue word "head" is semantically related to "person" (e.g., head of a person) and "president" (e.g., head of state), but the response format was rejected. Based on Trial \#2, the issue was likely the formatting of the TARGETS line or extra text. I will select a single{-}word clue that is not in any list (e.g., "mind" for mind, army, middle) and strictly output only the two required lines: "CLUE: <WORD>" followed immediately by "TARGETS: <word1>, <word2>, ..." with no additional explanation, ensuring the clue is not one of the target words itself.} \\
\\ 
\\ 
\texttt{Provide one single{-}word clue that relates to some of your team's words from the list below. You can choose to how many and to which words your clue relates to.} \\
\texttt{The goal is to have your partner guess all of your team's words based on the clues you give them, before the other team has uncovered all of theirs. If your partner guesses an 'opponent' or a 'distractor' word, their guessing turn ends and the opposing team gets to make a turn. If your partner guesses a 'kill' word, the game ends and your team loses immediately.} \\
\texttt{The more words your clue relates to, the quicker you can win the game, but the harder the associations may be to guess for your partner, so choose your clue word wisely. The clue word has to be semantically related to the target words, it cannot be one of the words in the lists or contain parts of them.} \\
\\ 
\texttt{Always give your single{-}word clue and your comma{-}separated list of related target words in the following format and make your answers as short as possible, never include any other text than is required in this form:} \\
\\ 
\texttt{CLUE: <WORD>} \\
\texttt{TARGETS: <TARGETS>} \\
\\ 
\texttt{Your team words are: way, cause, report, mind, kid, army, middle, president.} \\
\texttt{Your opponent's team words are: law, doctor, world, city, record, gold, fool.} \\
\texttt{Distractor words are: story, coffee, light, evening, company, ship.} \\
\texttt{The kill words are: hospital.} \\
            }
        }
    }
    & \\ \\

    \theutterance \stepcounter{utterance}  
    & \multicolumn{4}{p{0.6\linewidth}}{
        \cellcolor[rgb]{0.9,0.9,0.9}{
            \makecell[{{p{\linewidth}}}]{
                \texttt{\tiny{[P1$\rangle$GM]}}
                \texttt{CLUE: mind} \\
\texttt{TARGETS: mind, army, middle} \\
            }
        }
    }
    & & \\ \\

    \theutterance \stepcounter{utterance}  
    & & & \multicolumn{2}{p{0.3\linewidth}}{
        \cellcolor[rgb]{0.9,0.9,0.9}{
            \makecell[{{p{\linewidth}}}]{
                \texttt{\tiny{[GM$|$GM]}}
                {'player': 'Cluegiver', 'type': 'clue is morphologically related to word on the board', 'utterance': 'CLUE: mind\nTARGETS: mind, army, middle', 'message': "Your clue 'mind' is morphologically similar to the word mind, please choose another clue word.", 'clue': 'mind', 'similar_board_word': 'mind'}
            }
        }
    }
    & & \\ \\

    \theutterance \stepcounter{utterance}  
    & \multicolumn{4}{p{0.6\linewidth}}{
        \cellcolor[rgb]{0.95,0.95,0.95}{
            \makecell[{{p{\linewidth}}}]{
                \texttt{\tiny{[P2$\rightarrow$P1]}}
                \texttt{Post{-}episode Reflexion (not sent to GM):} \\
\texttt{CLUE: mind} \\
\texttt{TARGETS: mind, army, middle} \\
            }
        }
    }
    & & \\ \\

\end{supertabular}
}

\end{document}
